\section{Síntesis y ampliación}

\subsection{Repaso de los puntos clave}

Esta clase partió de la intersección entre la IA y la biología para presentar en profundidad las aplicaciones de vanguardia de la IA generativa en el diseño de proteínas; el contenido central puede resumirse en los siguientes niveles:

\begin{enumerate}
    \item \textbf{La biología es uno de los campos de aplicación con mayor impacto para la IA}: las proteínas, como «ejecutoras» de la biología, y los avances en la capacidad de diseñarlas influirán profundamente en el desarrollo de fármacos, la ingeniería de enzimas, la biología sintética y otros campos

    \item \textbf{La jerarquía de tres niveles secuencia-estructura-función de las proteínas} es el marco básico para entender y manipular las proteínas; el diseño de proteínas impulsado por IA busca invertir esta jerarquía, diseñando la secuencia de manera inversa a partir de la función

    \item \textbf{El modelo de difusión (Diffusion Model) puede trasladarse del dominio continuo al dominio discreto}: mediante estrategias de Masking o Mutation se define el proceso de adición de ruido sobre datos discretos, con lo que se establece un marco completo de difusión discreta

    \item \textbf{EvoDiff es un modelo de generación de secuencias de proteínas basado en difusión discreta}: se entrenó con aproximadamente 50 millones de secuencias a escala evolutiva y admite tanto la generación incondicional como la generación condicionada por función (Motif Scaffolding)

    \item \textbf{La validación experimental es el criterio final}: las proteínas generadas por EvoDiff ya se han expresado con éxito en el laboratorio, y se ha verificado tanto su estabilidad estructural como su actividad funcional (unión de iones de calcio)

    \item \textbf{AI for Biology es una ingeniería de sistemas que abarca múltiples escalas}: desde la molécula hasta la célula, el tejido y el individuo, se requiere la integración y colaboración de métodos de IA multimodal
\end{enumerate}

\subsection{El significado teórico de los modelos de difusión discreta}

Desde la perspectiva de la teoría del aprendizaje profundo, esta clase reveló una perspectiva unificadora profunda: el marco de Masking Diffusion discreto unifica matemáticamente los modelos de lenguaje autorregresivos y el Masked Language Model, y constituye un paradigma de modelado generativo más general. Esta contribución teórica no solo tiene valor para el diseño de proteínas, sino que también resulta esclarecedora para tareas más amplias de generación de datos discretos (como el diseño de moléculas, la generación de código, etc.).

\subsection{Retos clave y direcciones futuras}

Aunque EvoDiff mostró resultados iniciales apasionantes, el campo todavía enfrenta retos importantes:

\begin{itemize}
    \item \textbf{Escalabilidad de la validación funcional}: la validación experimental actual todavía se limita al nivel de casos individuales; cómo escalarla al nivel de sistemas es un cuello de botella clave
    \item \textbf{Modelado conjunto de secuencia, estructura y función}: cómo fusionar de manera eficaz la información de distintas modalidades y, al mismo tiempo, manejar las enormes diferencias en la escala de los datos
    \item \textbf{Interfaces de diseño impulsadas por lenguaje natural}: lograr un sistema de extremo a extremo que «describa los requisitos funcionales en lenguaje natural $\rightarrow$ genere automáticamente la proteína que cumpla esos requisitos»
    \item \textbf{Optimización experimental en ciclo cerrado}: establecer un ciclo iterativo cerrado de diseño con IA $\rightarrow$ prueba experimental $\rightarrow$ optimización por retroalimentación (active learning)
    \item \textbf{Integración entre escalas}: modelado de extremo a extremo desde el diseño molecular hasta la respuesta celular y los efectos a nivel de tejido o paciente
\end{itemize}

\subsection{Lecturas adicionales}

\begin{itemize}
    \item \textbf{Artículo de EvoDiff}：Alamdari et al., ``Protein generation with evolutionary diffusion: sequence is all you need'', \textit{Nature Biotechnology}, 2023
    \item \textbf{AlphaFold}：Jumper et al., ``Highly accurate protein structure prediction with AlphaFold'', \textit{Nature}, 2021
    \item \textbf{RFdiffusion} (método de difusión estructural)：Watson et al., ``De novo design of protein structure and function with RFdiffusion'', \textit{Nature}, 2023
    \item \textbf{Modelo de lenguaje de proteínas ESM}：Lin et al., ``Evolutionary-scale prediction of atomic-level protein structure with a language model'', \textit{Science}, 2023
    \item \textbf{Revisión sobre el modelado conjunto de secuencia y estructura}：Wang et al., artículo de revisión del equipo de MSR (mencionado en la clase)
    \item \textbf{Página del curso MIT 6.S191}：\url{http://introtodeeplearning.com/}
    \item \textbf{Equipo de MSR BioML}：\url{https://www.microsoft.com/en-us/research/}
\end{itemize}

%% --- Fin del contenido --- %%

\end{document}
